% SC26 reproducibility appendix — template from
% https://github.com/jennfshr/sc26-repro/tree/main/for-paper-authors
\documentclass[conference]{IEEEtran}

\usepackage{cite}
\usepackage{amsmath,amssymb,amsfonts}
\usepackage{algorithmic}
\usepackage{graphicx}
\usepackage{textcomp}
\usepackage{xcolor}
\usepackage{booktabs}
\usepackage{tagging}

\usepackage{sc26repro}
\usepackage{url}
\emergencystretch=3em  % forgive long \texttt tokens in 2-column
\sloppy

% Uncomment/comment the following usetag lines
% if you would like to get an explanation or an example.
% Don't forget to comment them for the final submission.
%\usetag{explanation}
%\usetag{example}

\begin{document}

\twocolumn[%
{\begin{center}
\Huge
Appendix: Artifact Description/Artifact Evaluation
\end{center}}
]

%%%%%%%%%%%%%%%%%%%%%%%%%%%%%%%%%%%%%%%%%%%%%%%%%%%%%
%  AD Appendix
%%%%%%%%%%%%%%%%%%%%%%%%%%%%%%%%%%%%%%%%%%%%%%%%%%%%%

\appendixAD

\section{Overview of Contributions and Artifacts}

\subsection{Paper's Main Contributions}

\begin{description}
\item[$C_1$] \textbf{GPU-accelerated I/O library interposition.} A VOL
  connector and filter for HDF5 (and an ADIOS2 operator) that
  transparently redirect scientific I/O calls onto a GPU-resident
  pipeline, keeping simulation data on the device from production
  through compression and storage without source-code changes to the
  application.
\item[$C_2$] \textbf{GPU-aware compression pipeline synthesis.} A
  lightweight neural-network selector driven by a Compression
  Characteristic Model that extracts eight per-chunk statistical
  descriptors and jointly picks the codec (among eight nvCOMP
  backends), byte-shuffle mode, quantization level, and error bound.
  Multi-level online refinement (cost-model scoring, SGD correction,
  top-$K$ exploration) adapts the selector to unseen workloads at
  runtime.
\item[$C_3$] \textbf{Resource-aware dynamic data organization.} A
  placement engine that routes compressed chunks across the HBM /
  DRAM / NVMe / HDD hierarchy using current tier capacity and
  bandwidth, asynchronously flushing to slower tiers during compute
  phases so data movement overlaps with productive work.
\end{description}


\subsection{Computational Artifacts}

\begin{description}
\item[$A_1$] NeuroPress core library + HDF5 VOL connector.
  Public source repository:
  \url{https://github.com/grc-iit/NeuroPress}.
\item[$A_2$] Neural-network training pipeline and synthetic
  benchmark. Same repository as $A_1$, under
  \texttt{neural\_net/} and \texttt{synthetic/}.
\item[$A_3$] Jarvis-CD pipelines for five scientific workloads
  (VPIC, Nyx, LAMMPS, WarpX, ResNet/ViT AI checkpoints),
  under \texttt{gpucompress\_pkgs/} in the $A_1$ repository.
  Built on the Jarvis-CD orchestration framework:
  \url{https://github.com/grc-iit/jarvis-cd}.
\item[$A_4$] Paper-evaluation harness and plot-generation scripts,
  under \texttt{evaluations/} and \texttt{analysis/} in the $A_1$
  repository.
\end{description}

\begin{center}
\begin{tabular}{rll}
\toprule
Artifact ID  &  Contributions &  Related \\
             &  Supported     &  Paper Elements \\
\midrule
$A_1$   &  $C_1, C_2, C_3$ & Fig.~3; Tab.~I \\
\midrule
$A_2$   &  $C_2$           & Fig.~4; \S IV.A \\
\midrule
$A_3$   &  $C_1, C_2, C_3$ & Fig.~5; Fig.~6; Fig.~8; Fig.~9;
                             Tab.~III \\
\midrule
$A_4$   &  $C_1, C_2, C_3$ & Figs.~4, 5, 6, 8, 9 (plotting) \\
\bottomrule
\end{tabular}
\end{center}


\subsection{Scope of Reproducibility}

One experiment in the paper --- the GPUDirect Storage (GDS)
bandwidth study in Fig.~7 --- was collected on an
NVIDIA-internal H100 + enterprise-NVMe testbed during an
author's internship and cannot be re-executed by external
reviewers. GDS requires a narrow stack alignment (GPU
firmware, NVMe firmware, kernel GDS module, file-system
driver, and matched CUDA/driver versions) that very few sites
have available; the hardware and platform access, not the
source code, are the limiting factor. Because the underlying
measurements are NVIDIA-internal, neither the raw CSVs nor a
dedicated plotter for Fig.~7 ship with this artifact; the
bars in the published Fig.~7 are a one-time archival result
and are out of scope for external reproduction.

All other results are fully reproducible on commodity A100
nodes (the NCSA Delta \texttt{gpuA100x4} partition was used
for the paper) through artifacts $A_1$--$A_4$ without any
NVIDIA-internal infrastructure:
Fig.~4 (NN selector accuracy and SHAP feature importance,
$C_2$) from $A_2$;
Fig.~5 (per-chunk cost breakdown pie charts, $C_2$) from
$A_3$'s figure\_5 pipelines (release-mode runs with
stage-timing instrumentation);
Fig.~6 (VPIC threshold-sensitivity sweep, $C_2$) from $A_3$;
Fig.~8 (cross-workload regret and MAPE convergence over five
workloads, $C_2$) from $A_3$; and
Fig.~9 (end-to-end wall-clock reduction of Baseline vs.
nvCOMP vs. NeuroPress configurations, $C_1$ and $C_3$) from
$A_3$. Tables~II and~III (cost-model weights and loss
configurations) are author-specified configuration tables,
not generated outputs.


%%%%%%%%%%%%%%%%%%%%%%%%%%%%%%%%%%%%%%%%%%%%%%%%%%%%%%%%
\section{Artifact Identification}
%%%%%%%%%%%%%%%%%%%%%%%%%%%%%%%%%%%%%%%%%%%%%%%%%%%%%%%%


\newartifact

\artrel
$A_1$ is the NeuroPress core library and substantiates all
three contributions; it is the implementation of the system.
For $C_1$ (transparent GPU-resident I/O path) it provides
\texttt{libgpucompress.so} (core runtime) and
\texttt{libH5VLgpucompress.so} (HDF5 VOL connector), which
together intercept standard HDF5 calls and route buffers
through a GPU-resident compression pipeline. For $C_2$
(GPU-aware compression pipeline synthesis), the library
implements the four blocks of the selector shown in Fig.~3:
the Compression Characteristic Model (per-chunk
statistical-descriptor extraction), Compression Selection (NN
forward pass over eight candidate pipelines, cost-model
re-ranking, online SGD correction, and top-$K$ exploration),
the Preprocess Factory (BitShuffle, Quantize), and the
Compressor Factory (ZLIB, LZMA, SZ3, and the eight nvCOMP
backends). The offline training that produces the weights
consumed by Compression Selection lives in $A_2$. For $C_3$
(resource-aware dynamic data organization), the library
contains the tier-aware placement engine --- the Operator
Offloader and Data Transfer Engine in Fig.~3 --- that routes
compressed chunks across the HBM/DRAM/NVMe/HDD hierarchy and
overlaps flushes with compute.

\artexp
Running the VOL-enabled unit tests (\texttt{ctest --test-dir
build}) should complete with all targets passing, with every
lossless round-trip returning bit-identical data and every
device-pointer write path bypassing the host bounce buffer
(verified via the xfer-audit counters in
\texttt{test\_vol\_xfer\_audit}). This substantiates $C_1$
(transparent GPU-resident I/O path) and confirms that the $C_3$
placement engine routes data through the full tiered-storage
path without corruption.

\arttime
The full \emph{build + test} cycle for $A_1$ completes in
under \textbf{10 minutes} end-to-end on an A100 node. Measured
breakdown (verified on Delta's \texttt{dt-login02}, 128 cores;
compute-node variance bounded by core count):
\texttt{./scripts/install\_deps.sh} (fetch pre-built nvcomp
5.1 archive from NVIDIA, then download, configure, build, and
install HDF5 2.0.0 with
\texttt{-DHDF5\_ENABLE\_PARALLEL=ON}) --- \textbf{1 min 36 s};
configure and build NeuroPress via
\texttt{cmake -B build \&\& cmake --build build -j\$(nproc)}
(library, 90 test binaries, MPI benchmarks)
--- \textbf{41 s};
\texttt{ctest --test-dir build --output-on-failure} (71 unit
/ VOL / regression / perf / nn tests) --- \textbf{$\le$ 2
minutes} on A100. Total: approximately \textbf{4 minutes on a
well-provisioned login node, $\le$ 10 minutes on a compute
node with fewer cores or slower links}. Analysis (inspect the
\texttt{ctest} summary and the per-test log under
\texttt{build/Testing/}) 2 minutes. No workload execution at
this stage --- workload runs are covered by $A_3$.

\artin

\artinpart{Hardware}
x86\_64 host with an NVIDIA GPU of compute capability $\geq$
7.0. Tested configurations: A100-40GB on NCSA Delta, RTX 4070 on
local development, H100 for the GDS evaluation. At least 32~GB
host RAM and 20~GB of free disk for the build tree.

\artinpart{Software}
Ubuntu 22.04 or RHEL 9 equivalent; CUDA 12.x; nvCOMP 5.x;
HDF5 1.14+ built with parallel support
(\texttt{-DHDF5\_ENABLE\_PARALLEL=ON}); OpenMPI 4.x or MPICH
4.x; GCC 12 (NVCC 12.8 has an ICE with libstdc++ 13); CMake
3.24+. The full dependency manifest is given in
\texttt{docker/Dockerfile} and
\texttt{cmake/CoreLibrary.cmake}.

\artinpart{Datasets / Inputs}
No external dataset is required for $A_1$. The synthetic
benchmark (\texttt{build/synthetic\_generator})
generates float32 chunks with controlled statistical properties
for unit validation.

\artinpart{Installation and Deployment}
\begin{verbatim}
git clone https://github.com/grc-iit/NeuroPress
cd NeuroPress
./scripts/install_deps.sh
cmake -B build
cmake --build build -j$(nproc)
ctest --test-dir build --output-on-failure
\end{verbatim}
\texttt{scripts/install\_deps.sh} stages HDF5-parallel 2.0.0
and nvcomp 5.1 under \texttt{.deps/} following the exact
recipe used by \texttt{docker/Dockerfile}; the Apptainer path
(used by $A_3$'s Jarvis pipelines) stages the same pair under
\texttt{/opt/deps/} inside the container. Bare-metal and
containerised builds therefore consume an identical dependency
tree. After the script runs, \texttt{cmake -B build}
auto-detects \texttt{HDF5\_ROOT} and \texttt{NVCOMP\_PREFIX}
from \texttt{.deps/}; no \texttt{-D} flags are required for
the default build. On non-Delta systems,
\texttt{CMAKE\_CUDA\_ARCHITECTURES} may need
\texttt{-DCMAKE\_CUDA\_ARCHITECTURES=<cc>} for a specific GPU
(default is \texttt{80} for A100).

\artcomp
Unit-level workflow: $T_1$ generates synthetic chunks, $T_2$
compresses and decompresses each chunk under every codec, and
$T_3$ verifies bitwise equality and records timing. Parameters:
4000 chunks, 32 (algorithm, shuffle, quant) configurations, seed
42.

\artout
\texttt{results/synthetic\_results.csv} (128k rows). A successful
build produces all three shared objects under \texttt{build/} and
all ctest targets pass.


\newartifact

\artrel
$A_2$ substantiates $C_2$: the neural compression selector. It
produces the packed-weight file
(\texttt{neural\_net/weights/model.nnwt}) consumed at runtime by
the VOL connector shipped in $A_1$.

Training the 8-input / 5-layer / 8-output fully-connected
network (13{,}312 parameters) on 128{,}000 synthetic rows
(2{,}000 chunks $\times$ 64 configurations, 80/20 by-chunk
split, 5-fold CV, Adam lr=$3\times10^{-4}$, weight
decay=$10^{-5}$, batch 512, up to 150 epochs with early
stopping) should converge to MAPE in the 10--23\% range
across all eight predicted statistics, reproducing Fig.~4(b)
of the paper; the SHAP feature-importance heatmap over the
eight inputs should reproduce Fig.~4(a), with compression
library, first derivative, and quantization level dominating
the predictions. Together these substantiate $C_2$ (GPU-aware
compression pipeline synthesis). The trained weights
(\texttt{neural\_net/weights/model.nnwt}) are consumed at
runtime by $A_3$'s figure\_5 pipelines to produce the
per-chunk stage-timing pie charts of Fig.~5, and by the
figure\_6/figure\_8 pipelines for online refinement (SGD +
top-$K$ exploration).

\arttime
Setup (Python environment and PyTorch) 10 minutes. Execution:
synthetic benchmark to generate the 128k-row training CSV
(\texttt{build/synthetic\_generator}) 15--20
minutes; training with five-fold cross-validation
(\texttt{neural\_net/training/train.py}) 25--35 minutes on one A100;
weight export (\texttt{neural\_net/export/export\_weights.py}) under 1 minute.
Analysis (regenerate the Fig.~4 accuracy and SHAP panels from
\texttt{results/}) 5 minutes. Total wall time: approximately
one hour on a single GPU.

\artin

\artinpart{Hardware}
Any CUDA-capable GPU with $\geq$ 8~GB VRAM. Training is not
latency-sensitive; a consumer GPU is sufficient.

\artinpart{Software}
Python 3.10+, PyTorch 2.2, torchvision, pyyaml, xgboost, shap,
matplotlib, seaborn, scikit-learn, pandas, numpy. The full
installation command is listed in
\texttt{docs/reproducability.md} under ``Prerequisites''.

\artinpart{Datasets / Inputs}
Generated in-tree. The synthetic benchmark produces the 128k-row
training CSV; no external download is required.

\artinpart{Installation and Deployment}
\begin{verbatim}
pip install torch torchvision pyyaml xgboost shap \
    matplotlib seaborn scikit-learn pandas numpy
build/synthetic_generator \
    --num-chunks 4000 \
    --output results/synthetic_results.csv \
    --seed 42
python neural_net/training/train.py \
    --csv results/synthetic_results.csv
# train.py writes neural_net/weights/model.pt
# The separate export step produces the .nnwt binary
# format that libgpucompress.so loads at runtime:
python neural_net/export/export_weights.py \
    --input  neural_net/weights/model.pt \
    --output neural_net/weights/model.nnwt
\end{verbatim}

Workflow:
$T_1$ synthetic data generation (2{,}000 float32 chunks at
four sizes: 64~KB, 256~KB, 1~MB, 4~MB, each tested under 64
configurations = 8 algorithms $\times$ 2 shuffle modes
$\times$ 4 quantization levels);
$T_2$ five-fold cross-validation
(\texttt{neural\_net/training/cross\_validate.py}) producing
the per-output MAPE bars of Fig.~4(b);
$T_3$ SHAP analysis on held-out chunks producing Fig.~4(a)'s
feature-importance heatmap;
$T_4$ weight export
(\texttt{neural\_net/export/export\_weights.py}) to the packed
\texttt{.nnwt} format consumed by $A_1$. Seeds are fixed to
42 throughout. Fig.~5 (per-chunk cost breakdown) is produced
by $A_3$'s figure\_5 pipelines at workload scale, not by a
microbenchmark here.

\artout
\texttt{neural\_net/weights/model.nnwt} (13\,312 parameters);
\texttt{results/training\_metrics.csv} (per-epoch loss and MAPE);
SHAP figures under \texttt{results/shap/}. The full procedure is
documented in \texttt{docs/reproducability.md}.


\newartifact

\artrel
$A_3$ substantiates $C_1$, $C_2$, and $C_3$ jointly by driving
the full stack on realistic simulations. Each Jarvis package
builds an Apptainer image containing the application linked
against GPUCompress and runs it with both a \emph{baseline}
(unmodified I/O) and an \emph{adaptive} (VOL + NN selector)
configuration.

\artexp
The adaptive configuration (\texttt{NeuroPress + Tier +
Async}) should reproduce four falsifiable predictions:
(i) \textbf{stage-timing anatomy (Fig.~5, per workload)}: the
figure\_5 pipelines emit six-wedge write-path and three-wedge
read-path pie charts over one H5Dwrite/H5Dread. The Stats
Kernel, NN Inference, and Compress Choice wedges should each
take $<15\%$ of the total write time, with Compress Time and
I/O Time dominating the remainder; the read path should show
Decompress Time and I/O Time accounting for $>90\%$ of total
read time;
(ii) \textbf{threshold sweep (Fig.~6, VPIC)}: cost MAPE and
regret rise monotonically with both online-learning thresholds
$X_1$ (proportional correction) and $X_2 - X_1$ (exploration);
write and read bandwidth peak near $X_1 = 30$, $X_2 - X_1 =
20$, giving the accuracy--throughput sweet spot adopted for
the remaining experiments;
(iii) \textbf{online convergence (Fig.~8, all five workloads)}:
LAMMPS and the AI checkpoint workload start near 15$\times$
and 8$\times$ regret respectively and decay to near-zero
within 200--500 chunks; VPIC, Nyx, and WarpX stabilise near
zero regret from the first chunk;
(iv) \textbf{end-to-end wall-clock (Fig.~9)}: NeuroPress +
Tier + Async reduces wall-clock time by approximately 53\%
versus the uncompressed baseline and 24\% versus
\texttt{nvCOMP + tiering} on I/O-heavy workloads (VPIC,
LAMMPS), 7--12\% on compute-heavy workloads (Nyx, WarpX, AI),
with an additional 10--15\% reduction when relaxing the
quality constraint from $\epsilon_L$ to $\epsilon_H$.

\arttime
Per workload on a single A100, broken out by phase:
(i) one-time image build and dependency resolution
(\texttt{jarvis ppl load yaml}) 15--20 minutes cold, cached
thereafter; (ii) SLURM allocation 1--5 minutes depending on
queue; (iii) actual benchmark execution (\texttt{jarvis ppl
run}, which drives both the baseline and adaptive variants
end-to-end and writes the CSVs listed under \artout):
approximately 15 seconds for the smoke configuration, 3
minutes for the small configuration, and 1.5 hours at paper
scale (the grid sizes and timestep counts reported in \S V);
(iv) archiving the outputs from the compute-node scratch into
\texttt{\$HOME} approximately 1 minute. All five workloads at
paper scale fit in the 8-hour AE budget on a single A100
allocation.

\artin

\artinpart{Hardware}
One NVIDIA A100-40GB (the NCSA Delta \texttt{gpuA100x4}
partition was used for the paper). Multi-GPU runs require the
grid dimension to divide the rank count cleanly; per-workload
READMEs list the valid combinations.

\artinpart{Software}
Apptainer 1.2+; Jarvis-CD (\texttt{pip install jarvis-cd}); CUDA
12.x host driver; SLURM for scheduling. All other dependencies
are baked into the Apptainer image produced by
\texttt{gpucompress\_base} and the per-workload
\texttt{gpucompress\_*\_delta} packages.

\artinpart{Datasets / Inputs}
All inputs are deterministically regenerated from shipped decks
and patches; no external dataset download (beyond CIFAR-10 for
the AI workload) is required. The on-disk output volume
produced during a run, however, scales with the simulation and
is tens of gigabytes at paper scale. Per-workload inputs:
VPIC --- Harris-sheet magnetic-reconnection deck, generated
in-tree at \texttt{benchmarks/vpic-kokkos/}; Nyx --- Sedov
blast wave ($256^3$) and Santa Barbara cosmology ($32^3$)
input decks, shipped in \texttt{benchmarks/nyx/patches/};
LAMMPS --- Lennard-Jones melt, generated in-tree; WarpX ---
laser--plasma acceleration input file, shipped in
\texttt{benchmarks/warpx/}; AI checkpointing --- CIFAR-10,
downloaded by \texttt{torchvision} on first run
(\textasciitilde 170~MB).

\artinpart{Installation and Deployment}
The Jarvis Path-B workflow replaces the earlier hand-written
shell scripts. Pipeline YAMLs live in two places: the
baseline pair (\texttt{<wl>\_baseline.yaml},
\texttt{<wl>\_single\_node.yaml}) under
\texttt{gpucompress\_pkgs/pipelines/} for $A_2/A_3$, and the
figure-specific YAMLs (trace-mode for Fig.~6,
release-mode-per-policy for Fig.~8) under
\texttt{evaluations/figure\_6/} and
\texttt{evaluations/figure\_8/}. Register the repo once, then
load whichever YAML matches the target figure:
\begin{verbatim}
# After the A_1 clone, export the repo root so every
# subsequent command is path-independent:
cd NeuroPress && export NEUROPRESS=$PWD
cd $NEUROPRESS
jarvis repo add gpucompress_pkgs

# Fig. 5 -- per-chunk cost breakdown (pie anatomy)
# Supported workloads for Fig. 5: VPIC, Nyx, WarpX
# (all three route through generic_benchmark which emits
#  stats_ms / nn_ms / comp_ms / read_ms / decomp_ms).
# LAMMPS's in-situ fix_gpucompress_kokkos does NOT currently
# emit those columns, so LAMMPS is not included in Fig. 5.
jarvis ppl load yaml \
  evaluations/figure_5/pipeline_<workload>.yaml

# Fig. 6 -- trace-mode threshold sweep (all four workloads)
jarvis ppl load yaml \
  evaluations/figure_6/pipeline_<workload>.yaml

# Fig. 8 -- release-mode per-policy (balanced/ratio)
jarvis ppl load yaml \
  evaluations/figure_8/pipeline_<workload>.yaml
\end{verbatim}
The first \texttt{jarvis ppl load yaml} invocation builds the
Apptainer image (15--20 minutes cold, cached thereafter). On
Delta's \texttt{pam\_slurm\_adopt}-gated nodes, the
per-pipeline hostfile under
\texttt{\$HOME/.ppi-jarvis/shared/} must reference the
currently-allocated compute node, otherwise \texttt{jarvis
ppl run} fails with \emph{no active jobs on this node}.
Refresh it immediately after \texttt{salloc}:
\begin{verbatim}
# On the GPU compute node, before each jarvis ppl run.
# <pipeline-dir> is the value of the YAML's "name:" field, e.g.
#   Fig. 5 : gpucompress_figure_5_<workload>
#   Fig. 6 : gpucompress_<workload>_delta_single_node
#   Fig. 8 : gpucompress_<workload>_delta_single_node
hostname > $HOME/hostfile_single.txt
cp $HOME/hostfile_single.txt \
   $HOME/.ppi-jarvis/shared/<pipeline-dir>/hostfile
jarvis ppl run
\end{verbatim}
Each subsequent \texttt{jarvis ppl run} executes the baseline
and adaptive runs back-to-back inside the container. Each
\texttt{gpucompress\_<workload>\_delta/README.md} gives the
full four-phase (clean, build, allocate, run-and-archive)
recipe; the figure-specific READMEs under
\texttt{evaluations/figure\_6/} and
\texttt{evaluations/figure\_8/} cover the trace-mode and
per-policy flows respectively.

\artcomp
For each workload $W$ (VPIC, Nyx, LAMMPS, WarpX, AI-ckpt):
$T_1(W)$ build the image; $T_2(W)$ run the baseline
(\texttt{\_baseline.yaml}); $T_3(W)$ run the adaptive variant
with \texttt{phase=nn-rl} and \texttt{policy=balanced};
$T_4(W)$ outputs land under the YAML-configurable
\texttt{results\_dir}. The figure-specific pipelines use
dedicated subdirectories:
\texttt{\$HOME/GPUCompress/tmp/figure\_6\_<wl>\_trace/} for
Fig.~6 trace-mode runs and
\texttt{\$HOME/GPUCompress/tmp/figure\_8\_<wl>\_<policy>/}
(with \texttt{<policy>} $\in$ \{\texttt{balanced},
\texttt{ratio}\}) for Fig.~8 per-policy runs. Baseline
\texttt{gpucompress\_pkgs/pipelines/} runs default to
\texttt{\$HOME/GPUCompress/tmp/<workload>/}. No post-run
archiving step is required. Key sweep parameters: three
policies (\texttt{speed}, \texttt{balanced}, \texttt{ratio});
error bounds $\in \{0.0, 10^{-3}, 10^{-2}\}$ for lossy runs;
grid sizes as documented per README.

\artout
Per run, under the relevant
\path{$HOME/GPUCompress/tmp/<outdir>/} (see \artcomp{} for
the three \texttt{<outdir>} conventions), the following files
are produced:
\path{benchmark_*_timesteps.csv} (per-snapshot write/read time
and ratio);
\path{benchmark_*_timestep_chunks.csv} (per-chunk features and
decisions);
\path{benchmark_*_ranking.csv} (Kendall $\tau$ and regret);
\path{gpucompress_trace.csv} (per-chunk $\times$ per-config
trace, 24-column CSV, produced only by the Fig.~6 trace-mode
pipelines);
and
\path{gpucompress_vol_summary.txt} (lifetime I/O totals).
These CSVs are the direct inputs to $A_4$.


\newartifact

$A_4$ substantiates the published evidence for $C_1$, $C_2$,
and $C_3$ by turning the raw CSVs produced by $A_2$ and $A_3$
into the headline figures that back each claim:
Fig.~4 (NN selector MAPE across outputs and SHAP feature
importance --- $C_2$) from $A_2$'s training metrics;
Fig.~5 (per-chunk cost breakdown pie charts on the write
and read paths --- $C_2$) from $A_3$'s figure\_5 pipelines'
\texttt{benchmark\_<dataset>\_timesteps.csv} (release mode
with stage-timing instrumentation);
Fig.~6 (VPIC threshold-sensitivity heatmaps for regret, MAPE,
write bandwidth, and read bandwidth --- $C_2$) from $A_3$;
Fig.~8 (cross-workload regret and cost MAPE convergence over
the five unseen workloads --- $C_2$) from $A_3$;
Fig.~9 (total wall-clock time across six configurations ---
Baseline, nvCOMP, nvCOMP+tiering, NeuroPress+tiering,
NeuroPress+tiering+async, and loss-config sweep --- on all
five workloads, $C_1$ and $C_3$) from $A_3$. Without $A_4$ the
reviewer would have per-run CSVs but no mechanical path from
those CSVs back to the figures in the paper.

Running the plotters over a directory of archived run outputs
reproduces:
Fig.~4 (NN accuracy + SHAP, from $A_2$'s training metrics);
Fig.~5 (per-chunk cost breakdown pie charts, from $A_3$'s
figure\_5 stage-timing CSVs);
Fig.~6 (four $7 \times 7$ VPIC threshold-sweep heatmaps);
and Fig.~8 and Fig.~9 from the per-workload CSVs. Each plot
regenerated from the shipped archive should be
pixel-comparable to the corresponding plot in the submitted
PDF. Fig.~7 (GPUDirect Storage bandwidth) is NVIDIA-internal
and out of scope for $A_4$ (see \S1.3).

Plot regeneration from pre-collected CSVs: approximately 10
minutes end-to-end on a laptop (Fig.~4 under 1 minute, Fig.~5
under 1 minute, Fig.~6 $\sim$3 minutes for the four heatmap
panels, Fig.~8 $\sim$2 minutes, Fig.~9 $\sim$2 minutes for the
stacked-bar panel across five workloads). No GPU is required
at this stage.

\artin

\artinpart{Hardware}
Any x86\_64 workstation with $\geq$ 8~GB RAM.

\artinpart{Software}
Python 3.10+; matplotlib, seaborn, pandas, numpy. No CUDA
required.

\artinpart{Datasets / Inputs}
The per-figure result directories produced by $A_3$ under
\texttt{\$HOME/GPUCompress/tmp/}:
\texttt{figure\_5\_<wl>/} (stage-timing CSVs for Fig.~5),
\texttt{figure\_6\_<wl>\_trace/} (per-chunk $\times$
per-config trace for Fig.~6 and Fig.~9), and
\texttt{figure\_8\_<wl>\_<policy>/} (release-mode per-policy
CSVs for Fig.~8). WarpX's Fig.~6 trace nests one level deeper
under \texttt{vol\_nn-rl+exp50/}; \texttt{plot\_trace.py} and
\texttt{throughput.py} both auto-detect this case.

\artinpart{Installation and Deployment}
Each figure has a dedicated plotter committed under
\texttt{evaluations/}, \texttt{analysis/}, or
\texttt{neural\_net/}:
\begin{verbatim}
# Fig. 4 (NN MAPE + SHAP panels)
neural_net/training/cross_validate.py
# Fig. 5 (pie anatomy of write/read path timing)
evaluations/figure_5/plot_anatomy.py
# Fig. 6 (threshold sweep + MAPE/regret plots)
analysis/plot_trace.py
# Fig. 7 is NVIDIA-internal (see Section 1.3) and not included.
# Fig. 8
evaluations/figure_8/plot.py
# Fig. 9
evaluations/figure_8/throughput.py
\end{verbatim}
All plotters take a directory of archived run outputs and
write PNG/PDF plus a numerical summary CSV alongside.

Each plotter is driven independently:
$T_1$ \texttt{python neural\_net/training/cross\_validate.py
--csv results/synthetic\_results.csv}
$\rightarrow$ Fig.~4 (MAPE bars; SHAP heatmap is emitted by
\texttt{neural\_net/training/hyperparam\_study.py});
$T_2$ \texttt{python evaluations/figure\_5/plot\_anatomy.py
<figure\_5\_<workload>/>} $\rightarrow$ Fig.~5 (write- and
read-path pie anatomy, reads
\texttt{benchmark\_<dataset>\_timesteps.csv} emitted by the
figure\_5 pipeline);
$T_3$ \texttt{python analysis/plot\_trace.py
<figure\_6\_<workload>\_trace/gpucompress\_trace.csv>}
$\rightarrow$ Fig.~6 (MAPE / regret line plots plus four
$40\times 40$ threshold heatmaps, all produced by a single
invocation);
$T_4$ \texttt{python evaluations/figure\_8/plot.py
--sc26-dir <tmp> --policies balanced ratio}
$\rightarrow$ Fig.~8 (regret + cost-MAPE convergence + per-metric
MAPE bar);
$T_5$ \texttt{python evaluations/figure\_8/throughput.py}
$\rightarrow$ Fig.~9 (wall-clock across eight configurations,
five workloads) via
\texttt{analysis/plot\_e2e.py:plot\_fig8}.

PNG and PDF figures under \texttt{evaluations/*/figures/} and
\texttt{analysis/figures/}, matching Figs.~4, 5, 6, 8, 9 in
the paper. A numerical summary CSV is written alongside each
figure to enable spot-checking against the values reported in
the text (e.g., 53\%, 24\%, 12\% reductions for Fig.~9;
$X_1=30$, $X_2-X_1=20$ sweet spot for Fig.~6).

\end{document}
